%=============================================================================
%  01_abstract.tex  -  Abstract + document conventions (roman-numbered)
%=============================================================================
\chapter*{Abstract}
\addcontentsline{toc}{chapter}{Abstract}
\markboth{Abstract}{}

\begin{tcolorbox}[enhanced,colback=uaMist,colframe=uaIndigo,boxrule=0pt,leftrule=4pt,
  arc=2pt,left=14pt,right=14pt,top=12pt,bottom=12pt]
\setlength{\parindent}{0pt}
\textbf{\large\color{uaIndigoD}Universal Analyst} is a self-hosted, LLM-driven
workspace that turns any uploaded dataset into an interrogable analytical
instrument. A user uploads a file --- CSV, JSON, Parquet, a PDF containing
tables, or a SQLite database --- and the system automatically \emph{profiles}
the data, runs \emph{exploratory data analysis}, and builds a
\emph{retrieval-augmented} context so that a \emph{local} multimodal large
language model (Gemma\,4, served through Ollama) can answer questions that are
grounded in the actual data rather than in the model's parametric memory.
\end{tcolorbox}

\vspace{6pt}
The platform goes beyond conversational question answering. Through a
router-driven \textbf{skill system}, the same session can request a
model-planned and server-\emph{validated} interactive dashboard, run a
lightweight \textbf{predictive-modeling} pipeline built on scikit-learn, perform
guided data \textbf{cleaning}, or export a branded, professional \textbf{PDF
report} containing the profile, quality signals, generated charts, and the full
conversation transcript. A central design commitment --- the \emph{no-fabrication
rule} --- ensures the interface never substitutes placeholder analysis for real
model output; when the model or the runtime fails, the system reports precisely
what happened.

This document describes the complete system: its motivation and product
philosophy, the end-to-end architecture, and each subsystem in depth ---
ingestion, profiling, chunking, the ChromaDB retrieval layer, the streaming
Ollama client, the six analytical skills, the Plotly chart-validation engine,
the scikit-learn prediction pipeline, the ReportLab PDF generator, the React
front end, the HTTP API surface, containerized deployment, testing, and
security posture. It closes with the team's six-way work split, a roadmap, and
appendices covering configuration and the event protocol.

\vspace{10pt}
\begin{center}
\begin{tikzpicture}
  \node[font=\itshape\large,text=uaSlate] {measured \;\textcolor{uaAmber}{\textbullet}\; grounded \;\textcolor{uaAmber}{\textbullet}\; sharp};
\end{tikzpicture}
\end{center}

\vfill
\subsection*{\color{uaIndigoD}Document Conventions}
\begin{itemize}
  \item \code{Inline monospace} denotes code identifiers, file names, endpoints,
        and configuration keys.
  \item Full code listings are drawn from the project's actual source and are
        line-numbered with an indigo margin rule.
  \item \textbf{Blue callouts} highlight design rationale; \textbf{amber
        callouts} flag constraints, caveats, and operational warnings.
  \item Figures render as vector diagrams; every figure and table is listed in
        the front matter.
\end{itemize}
